\section{R ile çözümlü uygulama}
\label{sec:r-b11}

\textbf{Soru.} Bölümdeki bağımsız grupların özetlerinden Welch testini;
yirmi öğrencinin fark puanlarının özetinden eşleştirilmiş testi hesaplayın.
Ham veriler verilmediğinden aşağıdaki çözümler özet istatistik formüllerini kullanır.

\begin{lstlisting}[style=prismR]
hacim <- c(30, 28)
ortalama <- c(78, 72)
standart_sapma <- c(8, 10)
katki <- standart_sapma^2 / hacim
standart_hata <- sqrt(sum(katki))
fark <- ortalama[1] - ortalama[2]
serbestlik <- sum(katki)^2 / sum(katki^2 / (hacim - 1))
t_degeri <- fark / standart_hata
c(t = t_degeri, sd = serbestlik,
  p = 2 * pt(abs(t_degeri), df = serbestlik,
             lower.tail = FALSE))
fark + c(-1, 1) * qt(0.975, df = serbestlik) * standart_hata
cift_sayisi <- 20
ortalama_fark <- 4.2
fark_sapmasi <- 5
es_se <- fark_sapmasi / sqrt(cift_sayisi)
es_t <- ortalama_fark / es_se
c(t = es_t, sd = cift_sayisi - 1,
  p = 2 * pt(abs(es_t), df = cift_sayisi - 1,
             lower.tail = FALSE),
  dz = ortalama_fark / fark_sapmasi)
ortalama_fark + c(-1, 1) *
  qt(0.975, df = cift_sayisi - 1) * es_se
\end{lstlisting}

\begin{outputguidebox}
\textbf{Çözüm ve kontrol değerleri.} Welch testinde fark 6 puan,
$t=2{,}5121$, serbestlik derecesi 51,7113, $p=0{,}01516$ ve yüzde 95
fark aralığı $[1{,}2066;10{,}7934]$'tür. Eşleştirilmiş testte
$t(19)=3{,}7566$, $p=0{,}001336$, aralık $[1{,}8599;6{,}5401]$ ve
$d_z=0{,}84$ bulunur.
\end{outputguidebox}

\textbf{Yorum.} Ham veriyle \texttt{t.test} kullanılırken bağımsız gruplar
için \texttt{var.equal = FALSE}, eşleşmiş son ve ön ölçüm vektörleri için
\texttt{paired = TRUE} seçilir. Eşleştirilmiş testte aynı
konumdaki iki değer aynı kişiye ait olmalıdır. Özetleri ham gözlem gibi bu
komutlara vermeyin; 51,7113 serbestlik derecesini de tam sayıya yuvarlamayın.

\textbf{Pekiştirme ve yanıt.} Fark tanımı ``ön--son'' yapılırsa eşleştirilmiş
$t$ ve $d_z$ işaret değiştirir; çift yönlü \pval-değeri değişmez ve güven
aralığı $[-6{,}5401;-1{,}8599]$ olur.